% Options for packages loaded elsewhere
\PassOptionsToPackage{unicode}{hyperref}
\PassOptionsToPackage{hyphens}{url}
%
\documentclass[
]{article}
\usepackage{amsmath,amssymb}
\usepackage{iftex}
\ifPDFTeX
  \usepackage[T1]{fontenc}
  \usepackage[utf8]{inputenc}
  \usepackage{textcomp} % provide euro and other symbols
\else % if luatex or xetex
  \usepackage{unicode-math} % this also loads fontspec
  \defaultfontfeatures{Scale=MatchLowercase}
  \defaultfontfeatures[\rmfamily]{Ligatures=TeX,Scale=1}
\fi
\usepackage{lmodern}
\ifPDFTeX\else
  % xetex/luatex font selection
\fi
% Use upquote if available, for straight quotes in verbatim environments
\IfFileExists{upquote.sty}{\usepackage{upquote}}{}
\IfFileExists{microtype.sty}{% use microtype if available
  \usepackage[]{microtype}
  \UseMicrotypeSet[protrusion]{basicmath} % disable protrusion for tt fonts
}{}
\makeatletter
\@ifundefined{KOMAClassName}{% if non-KOMA class
  \IfFileExists{parskip.sty}{%
    \usepackage{parskip}
  }{% else
    \setlength{\parindent}{0pt}
    \setlength{\parskip}{6pt plus 2pt minus 1pt}}
}{% if KOMA class
  \KOMAoptions{parskip=half}}
\makeatother
\usepackage{xcolor}
\usepackage[margin=1in]{geometry}
\usepackage{color}
\usepackage{fancyvrb}
\newcommand{\VerbBar}{|}
\newcommand{\VERB}{\Verb[commandchars=\\\{\}]}
\DefineVerbatimEnvironment{Highlighting}{Verbatim}{commandchars=\\\{\}}
% Add ',fontsize=\small' for more characters per line
\usepackage{framed}
\definecolor{shadecolor}{RGB}{248,248,248}
\newenvironment{Shaded}{\begin{snugshade}}{\end{snugshade}}
\newcommand{\AlertTok}[1]{\textcolor[rgb]{0.94,0.16,0.16}{#1}}
\newcommand{\AnnotationTok}[1]{\textcolor[rgb]{0.56,0.35,0.01}{\textbf{\textit{#1}}}}
\newcommand{\AttributeTok}[1]{\textcolor[rgb]{0.13,0.29,0.53}{#1}}
\newcommand{\BaseNTok}[1]{\textcolor[rgb]{0.00,0.00,0.81}{#1}}
\newcommand{\BuiltInTok}[1]{#1}
\newcommand{\CharTok}[1]{\textcolor[rgb]{0.31,0.60,0.02}{#1}}
\newcommand{\CommentTok}[1]{\textcolor[rgb]{0.56,0.35,0.01}{\textit{#1}}}
\newcommand{\CommentVarTok}[1]{\textcolor[rgb]{0.56,0.35,0.01}{\textbf{\textit{#1}}}}
\newcommand{\ConstantTok}[1]{\textcolor[rgb]{0.56,0.35,0.01}{#1}}
\newcommand{\ControlFlowTok}[1]{\textcolor[rgb]{0.13,0.29,0.53}{\textbf{#1}}}
\newcommand{\DataTypeTok}[1]{\textcolor[rgb]{0.13,0.29,0.53}{#1}}
\newcommand{\DecValTok}[1]{\textcolor[rgb]{0.00,0.00,0.81}{#1}}
\newcommand{\DocumentationTok}[1]{\textcolor[rgb]{0.56,0.35,0.01}{\textbf{\textit{#1}}}}
\newcommand{\ErrorTok}[1]{\textcolor[rgb]{0.64,0.00,0.00}{\textbf{#1}}}
\newcommand{\ExtensionTok}[1]{#1}
\newcommand{\FloatTok}[1]{\textcolor[rgb]{0.00,0.00,0.81}{#1}}
\newcommand{\FunctionTok}[1]{\textcolor[rgb]{0.13,0.29,0.53}{\textbf{#1}}}
\newcommand{\ImportTok}[1]{#1}
\newcommand{\InformationTok}[1]{\textcolor[rgb]{0.56,0.35,0.01}{\textbf{\textit{#1}}}}
\newcommand{\KeywordTok}[1]{\textcolor[rgb]{0.13,0.29,0.53}{\textbf{#1}}}
\newcommand{\NormalTok}[1]{#1}
\newcommand{\OperatorTok}[1]{\textcolor[rgb]{0.81,0.36,0.00}{\textbf{#1}}}
\newcommand{\OtherTok}[1]{\textcolor[rgb]{0.56,0.35,0.01}{#1}}
\newcommand{\PreprocessorTok}[1]{\textcolor[rgb]{0.56,0.35,0.01}{\textit{#1}}}
\newcommand{\RegionMarkerTok}[1]{#1}
\newcommand{\SpecialCharTok}[1]{\textcolor[rgb]{0.81,0.36,0.00}{\textbf{#1}}}
\newcommand{\SpecialStringTok}[1]{\textcolor[rgb]{0.31,0.60,0.02}{#1}}
\newcommand{\StringTok}[1]{\textcolor[rgb]{0.31,0.60,0.02}{#1}}
\newcommand{\VariableTok}[1]{\textcolor[rgb]{0.00,0.00,0.00}{#1}}
\newcommand{\VerbatimStringTok}[1]{\textcolor[rgb]{0.31,0.60,0.02}{#1}}
\newcommand{\WarningTok}[1]{\textcolor[rgb]{0.56,0.35,0.01}{\textbf{\textit{#1}}}}
\usepackage{graphicx}
\makeatletter
\def\maxwidth{\ifdim\Gin@nat@width>\linewidth\linewidth\else\Gin@nat@width\fi}
\def\maxheight{\ifdim\Gin@nat@height>\textheight\textheight\else\Gin@nat@height\fi}
\makeatother
% Scale images if necessary, so that they will not overflow the page
% margins by default, and it is still possible to overwrite the defaults
% using explicit options in \includegraphics[width, height, ...]{}
\setkeys{Gin}{width=\maxwidth,height=\maxheight,keepaspectratio}
% Set default figure placement to htbp
\makeatletter
\def\fps@figure{htbp}
\makeatother
\setlength{\emergencystretch}{3em} % prevent overfull lines
\providecommand{\tightlist}{%
  \setlength{\itemsep}{0pt}\setlength{\parskip}{0pt}}
\setcounter{secnumdepth}{-\maxdimen} % remove section numbering
\ifLuaTeX
  \usepackage{selnolig}  % disable illegal ligatures
\fi
\usepackage{bookmark}
\IfFileExists{xurl.sty}{\usepackage{xurl}}{} % add URL line breaks if available
\urlstyle{same}
\hypersetup{
  pdftitle={RNN and LSTM},
  pdfauthor={Kachun Huang},
  hidelinks,
  pdfcreator={LaTeX via pandoc}}

\title{RNN and LSTM}
\author{Kachun Huang}
\date{2025-04-29}

\begin{document}
\maketitle

\section{RNN and LSTM}\label{rnn-and-lstm}

\subsection{Import Library}\label{import-library}

\begin{Shaded}
\begin{Highlighting}[]
\FunctionTok{library}\NormalTok{(keras)}
\FunctionTok{library}\NormalTok{(tensorflow)}
\FunctionTok{library}\NormalTok{(reticulate)}
\end{Highlighting}
\end{Shaded}

\begin{verbatim}
## Warning: package 'reticulate' was built under R version 4.3.3
\end{verbatim}

\begin{Shaded}
\begin{Highlighting}[]
\FunctionTok{library}\NormalTok{(tidyverse)}
\end{Highlighting}
\end{Shaded}

\begin{verbatim}
## Warning: package 'purrr' was built under R version 4.3.3
\end{verbatim}

\begin{verbatim}
## Warning: package 'lubridate' was built under R version 4.3.3
\end{verbatim}

\begin{verbatim}
## -- Attaching core tidyverse packages ------------------------ tidyverse 2.0.0 --
## v dplyr     1.1.4     v readr     2.1.5
## v forcats   1.0.0     v stringr   1.5.1
## v ggplot2   3.5.1     v tibble    3.2.1
## v lubridate 1.9.4     v tidyr     1.3.1
## v purrr     1.0.4     
## -- Conflicts ------------------------------------------ tidyverse_conflicts() --
## x dplyr::filter() masks stats::filter()
## x dplyr::lag()    masks stats::lag()
## i Use the conflicted package (<http://conflicted.r-lib.org/>) to force all conflicts to become errors
\end{verbatim}

\begin{Shaded}
\begin{Highlighting}[]
\FunctionTok{library}\NormalTok{(tidyquant)}
\end{Highlighting}
\end{Shaded}

\begin{verbatim}
## Warning: package 'tidyquant' was built under R version 4.3.3
\end{verbatim}

\begin{verbatim}
## Registered S3 method overwritten by 'quantmod':
##   method            from
##   as.zoo.data.frame zoo
\end{verbatim}

\begin{verbatim}
## Warning: package 'xts' was built under R version 4.3.3
\end{verbatim}

\begin{verbatim}
## Warning: package 'zoo' was built under R version 4.3.3
\end{verbatim}

\begin{verbatim}
## Warning: package 'quantmod' was built under R version 4.3.3
\end{verbatim}

\begin{verbatim}
## Warning: package 'TTR' was built under R version 4.3.3
\end{verbatim}

\begin{verbatim}
## Warning: package 'PerformanceAnalytics' was built under R version 4.3.3
\end{verbatim}

\begin{verbatim}
## -- Attaching core tidyquant packages ----------------------- tidyquant 1.0.11 --
## v PerformanceAnalytics 2.0.8      v TTR                  0.24.4
## v quantmod             0.4.27     v xts                  0.14.1
## -- Conflicts ------------------------------------------ tidyquant_conflicts() --
## x zoo::as.Date()                 masks base::as.Date()
## x zoo::as.Date.numeric()         masks base::as.Date.numeric()
## x dplyr::filter()                masks stats::filter()
## x xts::first()                   masks dplyr::first()
## x dplyr::lag()                   masks stats::lag()
## x xts::last()                    masks dplyr::last()
## x PerformanceAnalytics::legend() masks graphics::legend()
## x quantmod::summary()            masks base::summary()
## i Use the conflicted package (<http://conflicted.r-lib.org/>) to force all conflicts to become errors
\end{verbatim}

\begin{Shaded}
\begin{Highlighting}[]
\FunctionTok{library}\NormalTok{(zoo)}
\FunctionTok{library}\NormalTok{(caret)}
\end{Highlighting}
\end{Shaded}

\begin{verbatim}
## Warning: package 'caret' was built under R version 4.3.3
\end{verbatim}

\begin{verbatim}
## Loading required package: lattice
\end{verbatim}

\begin{verbatim}
## Warning: package 'lattice' was built under R version 4.3.3
\end{verbatim}

\begin{verbatim}
## 
## Attaching package: 'caret'
## 
## The following object is masked from 'package:purrr':
## 
##     lift
## 
## The following object is masked from 'package:tensorflow':
## 
##     train
\end{verbatim}

\begin{Shaded}
\begin{Highlighting}[]
\FunctionTok{library}\NormalTok{(tidyquant)}
\FunctionTok{library}\NormalTok{(magrittr)}
\end{Highlighting}
\end{Shaded}

\begin{verbatim}
## 
## Attaching package: 'magrittr'
## 
## The following object is masked from 'package:purrr':
## 
##     set_names
## 
## The following object is masked from 'package:tidyr':
## 
##     extract
\end{verbatim}

\subsection{Data Pre-processing}\label{data-pre-processing}

Before building models, raw data often requires transformation. Common
pre-processing steps include:

\subsubsection{\texorpdfstring{\textbf{Normalization}}{Normalization}}\label{normalization}

To stabilize model learning, normalize features such as stock prices
using min-max scaling:

\begin{Shaded}
\begin{Highlighting}[]
\NormalTok{data}\SpecialCharTok{$}\NormalTok{min\_lagged }\OtherTok{\textless{}{-}} \FunctionTok{lag}\NormalTok{(data}\SpecialCharTok{$}\NormalTok{Low)}
\NormalTok{data}\SpecialCharTok{$}\NormalTok{max\_lagged }\OtherTok{\textless{}{-}} \FunctionTok{lag}\NormalTok{(data}\SpecialCharTok{$}\NormalTok{High)}
\NormalTok{data}\SpecialCharTok{$}\NormalTok{Close\_norm }\OtherTok{\textless{}{-}}\NormalTok{ (data}\SpecialCharTok{$}\NormalTok{Close }\SpecialCharTok{{-}}\NormalTok{ data}\SpecialCharTok{$}\NormalTok{min\_lagged) }\SpecialCharTok{/}\NormalTok{ (data}\SpecialCharTok{$}\NormalTok{max\_lagged }\SpecialCharTok{{-}}\NormalTok{ data}\SpecialCharTok{$}\NormalTok{min\_lagged)}
\end{Highlighting}
\end{Shaded}

\subsubsection{\texorpdfstring{\textbf{Handling Missing
Values}}{Handling Missing Values}}\label{handling-missing-values}

\begin{Shaded}
\begin{Highlighting}[]
\NormalTok{data }\OtherTok{\textless{}{-}} \FunctionTok{na.omit}\NormalTok{(data)}
\end{Highlighting}
\end{Shaded}

\subsubsection{Date Conversion
(Optional)}\label{date-conversion-optional}

\begin{Shaded}
\begin{Highlighting}[]
\NormalTok{data}\SpecialCharTok{$}\NormalTok{Date }\OtherTok{\textless{}{-}} \FunctionTok{as.Date}\NormalTok{(data}\SpecialCharTok{$}\NormalTok{Date)}
\end{Highlighting}
\end{Shaded}

\subsubsection{Train-Test Splitting}\label{train-test-splitting}

Split your data into training and test sets to evaluate your model's
generalization performance:

\begin{Shaded}
\begin{Highlighting}[]
\NormalTok{split\_index }\OtherTok{\textless{}{-}} \FunctionTok{floor}\NormalTok{(}\FloatTok{0.9} \SpecialCharTok{*} \FunctionTok{length}\NormalTok{(series))}
\NormalTok{train\_data }\OtherTok{\textless{}{-}}\NormalTok{ series[}\DecValTok{1}\SpecialCharTok{:}\NormalTok{split\_index]}
\NormalTok{test\_data }\OtherTok{\textless{}{-}}\NormalTok{ series[(split\_index }\SpecialCharTok{+} \DecValTok{1}\NormalTok{)}\SpecialCharTok{:}\FunctionTok{length}\NormalTok{(series)]}
\end{Highlighting}
\end{Shaded}

\begin{itemize}
\item
  \texttt{length(series)}~returns the total number of time steps in your
  dataset.
\item
  \texttt{floor(0.9\ *\ length(series))}~calculates an index at 90\% of
  the data.
\item
  \texttt{train\_data}~selects the first 90\% as the training set.
\item
  \texttt{test\_data}~selects the remaining 10\% as the test set.
\end{itemize}

\subsubsection{Creating Time Series
Datasets}\label{creating-time-series-datasets}

\begin{Shaded}
\begin{Highlighting}[]
\NormalTok{create\_dataset }\OtherTok{\textless{}{-}} \ControlFlowTok{function}\NormalTok{(series, }\AttributeTok{timesteps =} \DecValTok{10}\NormalTok{) \{}
\NormalTok{  x }\OtherTok{\textless{}{-}} \FunctionTok{list}\NormalTok{()}
\NormalTok{  y }\OtherTok{\textless{}{-}} \FunctionTok{c}\NormalTok{()}
  \ControlFlowTok{for}\NormalTok{ (i }\ControlFlowTok{in} \DecValTok{1}\SpecialCharTok{:}\NormalTok{(}\FunctionTok{length}\NormalTok{(series) }\SpecialCharTok{{-}}\NormalTok{ timesteps)) \{}
\NormalTok{    x[[i]] }\OtherTok{\textless{}{-}}\NormalTok{ series[i}\SpecialCharTok{:}\NormalTok{(i }\SpecialCharTok{+}\NormalTok{ timesteps }\SpecialCharTok{{-}} \DecValTok{1}\NormalTok{)]}
\NormalTok{    y }\OtherTok{\textless{}{-}} \FunctionTok{c}\NormalTok{(y, series[i }\SpecialCharTok{+}\NormalTok{ timesteps])}
\NormalTok{  \}}
\NormalTok{  x\_array }\OtherTok{\textless{}{-}} \FunctionTok{do.call}\NormalTok{(rbind, x)}
  \FunctionTok{return}\NormalTok{(}\FunctionTok{list}\NormalTok{(}\AttributeTok{x =}\NormalTok{ x\_array, }\AttributeTok{y =}\NormalTok{ y))}
\NormalTok{\}}
\end{Highlighting}
\end{Shaded}

This prepares the data in a format compatible with RNNs and LSTMs.

Then apply the~\texttt{create\_dataset()}~function:

\begin{Shaded}
\begin{Highlighting}[]
\NormalTok{train\_set }\OtherTok{\textless{}{-}} \FunctionTok{create\_dataset}\NormalTok{(train\_data, timesteps)}
\NormalTok{test\_set }\OtherTok{\textless{}{-}} \FunctionTok{create\_dataset}\NormalTok{(test\_data, timesteps)}
\end{Highlighting}
\end{Shaded}

Then split \texttt{train\_set} to \texttt{x\_train}, \texttt{y\_train}
and convert it to tensor

\begin{Shaded}
\begin{Highlighting}[]
\CommentTok{\#convert to array}
\NormalTok{x\_train\_array }\OtherTok{\textless{}{-}} \FunctionTok{array}\NormalTok{(}\FunctionTok{as.numeric}\NormalTok{(train\_set}\SpecialCharTok{$}\NormalTok{x), }\AttributeTok{dim =} \FunctionTok{c}\NormalTok{(}\FunctionTok{nrow}\NormalTok{(train\_set}\SpecialCharTok{$}\NormalTok{x), timesteps, }\DecValTok{1}\NormalTok{))}
\NormalTok{y\_train\_array }\OtherTok{\textless{}{-}} \FunctionTok{as.numeric}\NormalTok{(train\_set}\SpecialCharTok{$}\NormalTok{y)}

\CommentTok{\# Convert to TensorFlow tensors}
\NormalTok{x\_train }\OtherTok{\textless{}{-}}\NormalTok{ tf}\SpecialCharTok{$}\FunctionTok{convert\_to\_tensor}\NormalTok{(x\_train\_array, }\AttributeTok{dtype =} \StringTok{"float32"}\NormalTok{)}
\NormalTok{y\_train }\OtherTok{\textless{}{-}}\NormalTok{ tf}\SpecialCharTok{$}\FunctionTok{convert\_to\_tensor}\NormalTok{(y\_train\_array, }\AttributeTok{dtype =} \StringTok{"float32"}\NormalTok{)}
\end{Highlighting}
\end{Shaded}

\subsection{RNN Model Building}\label{rnn-model-building}

\begin{Shaded}
\begin{Highlighting}[]
\NormalTok{input\_rnn }\OtherTok{\textless{}{-}} \FunctionTok{layer\_input}\NormalTok{(}\AttributeTok{shape =} \FunctionTok{c}\NormalTok{(timesteps, }\DecValTok{1}\NormalTok{))}

\NormalTok{output\_rnn }\OtherTok{\textless{}{-}}\NormalTok{ input\_rnn }\SpecialCharTok{\%\textgreater{}\%}
  \FunctionTok{layer\_simple\_rnn}\NormalTok{(}\AttributeTok{units =} \DecValTok{32}\NormalTok{) }\SpecialCharTok{\%\textgreater{}\%}
  \FunctionTok{layer\_dense}\NormalTok{(}\AttributeTok{units =} \DecValTok{1}\NormalTok{)}

\NormalTok{model\_rnn }\OtherTok{\textless{}{-}} \FunctionTok{keras\_model}\NormalTok{(}\AttributeTok{inputs =}\NormalTok{ input\_rnn, }\AttributeTok{outputs =}\NormalTok{ output\_rnn)}
\end{Highlighting}
\end{Shaded}

\subsection{LSTM Model Building}\label{lstm-model-building}

\begin{Shaded}
\begin{Highlighting}[]
\NormalTok{input\_lstm }\OtherTok{\textless{}{-}} \FunctionTok{layer\_input}\NormalTok{(}\AttributeTok{shape =} \FunctionTok{c}\NormalTok{(timesteps, }\DecValTok{1}\NormalTok{))}

\NormalTok{output\_lstm }\OtherTok{\textless{}{-}}\NormalTok{ input\_lstm }\SpecialCharTok{\%\textgreater{}\%}
  \FunctionTok{layer\_lstm}\NormalTok{(}\AttributeTok{units =} \DecValTok{32}\NormalTok{) }\SpecialCharTok{\%\textgreater{}\%}
  \FunctionTok{layer\_dense}\NormalTok{(}\AttributeTok{units =} \DecValTok{1}\NormalTok{)}

\NormalTok{model\_lstm }\OtherTok{\textless{}{-}} \FunctionTok{keras\_model}\NormalTok{(}\AttributeTok{inputs =}\NormalTok{ input\_lstm, }\AttributeTok{outputs =}\NormalTok{ output\_lstm)}
\end{Highlighting}
\end{Shaded}

These models can be used for sequence-to-one prediction problems.

\subsection{\texorpdfstring{\textbf{Compilation and
Training}}{Compilation and Training}}\label{compilation-and-training}

\begin{Shaded}
\begin{Highlighting}[]
\NormalTok{model}\SpecialCharTok{$}\FunctionTok{compile}\NormalTok{(}
  \AttributeTok{optimizer =} \FunctionTok{optimizer\_adam}\NormalTok{(),}
  \AttributeTok{loss =} \StringTok{"mean\_squared\_error"}
\NormalTok{)}

\NormalTok{history }\OtherTok{\textless{}{-}}\NormalTok{ model}\SpecialCharTok{$}\FunctionTok{fit}\NormalTok{(}
  \AttributeTok{x =}\NormalTok{ x\_train,}
  \AttributeTok{y =}\NormalTok{ y\_train,}
  \AttributeTok{epochs =} \FunctionTok{r\_to\_py}\NormalTok{(}\DecValTok{100}\DataTypeTok{L}\NormalTok{),}
  \AttributeTok{batch\_size =} \FunctionTok{r\_to\_py}\NormalTok{(}\DecValTok{16}\DataTypeTok{L}\NormalTok{),}
  \AttributeTok{verbose =} \DecValTok{0}
\NormalTok{)}
\end{Highlighting}
\end{Shaded}

\texttt{r\_to\_py()}~ensures integers are converted correctly for
Python-based TensorFlow backend.

\subsection{Prediction}\label{prediction}

\begin{Shaded}
\begin{Highlighting}[]
\NormalTok{predict\_next }\OtherTok{\textless{}{-}} \ControlFlowTok{function}\NormalTok{(model, input\_seq\_array, }\AttributeTok{n\_steps =} \DecValTok{10}\NormalTok{) \{}
\NormalTok{  preds }\OtherTok{\textless{}{-}} \FunctionTok{c}\NormalTok{() }\CommentTok{\# empty vector to store prediction}
\NormalTok{  current\_input }\OtherTok{\textless{}{-}}\NormalTok{ input\_seq\_array }
  \ControlFlowTok{for}\NormalTok{ (i }\ControlFlowTok{in} \DecValTok{1}\SpecialCharTok{:}\NormalTok{n\_steps) \{}
\NormalTok{    pred }\OtherTok{\textless{}{-}}\NormalTok{ model}\SpecialCharTok{$}\FunctionTok{predict}\NormalTok{(current\_input)}
\NormalTok{    preds }\OtherTok{\textless{}{-}} \FunctionTok{c}\NormalTok{(preds, }\FunctionTok{as.numeric}\NormalTok{(pred))}
\NormalTok{    updated\_seq }\OtherTok{\textless{}{-}} \FunctionTok{c}\NormalTok{(}\FunctionTok{as.numeric}\NormalTok{(current\_input)[}\SpecialCharTok{{-}}\DecValTok{1}\NormalTok{], }\FunctionTok{as.numeric}\NormalTok{(pred))}
\NormalTok{    current\_input }\OtherTok{\textless{}{-}} \FunctionTok{array}\NormalTok{(updated\_seq, }\AttributeTok{dim =} \FunctionTok{c}\NormalTok{(}\DecValTok{1}\NormalTok{, timesteps, }\DecValTok{1}\NormalTok{))}
\NormalTok{    current\_input }\OtherTok{\textless{}{-}}\NormalTok{ tf}\SpecialCharTok{$}\FunctionTok{convert\_to\_tensor}\NormalTok{(current\_input, }\AttributeTok{dtype =} \StringTok{"float32"}\NormalTok{)}
\NormalTok{  \}}
  \FunctionTok{return}\NormalTok{(preds)}
\NormalTok{\}}
\end{Highlighting}
\end{Shaded}

This function generates~\texttt{n\_steps}~future predictions based on a
rolling window.

\subsubsection{Rescaling Predictions}\label{rescaling-predictions}

If you used normalization, reverse it using the original min and max
values:

\begin{Shaded}
\begin{Highlighting}[]
\NormalTok{rescaled\_value }\OtherTok{\textless{}{-}}\NormalTok{ norm\_value }\SpecialCharTok{*}\NormalTok{ (original\_max }\SpecialCharTok{{-}}\NormalTok{ original\_min) }\SpecialCharTok{+}\NormalTok{ original\_min}
\end{Highlighting}
\end{Shaded}

\subsection{Model Performance}\label{model-performance}

\begin{Shaded}
\begin{Highlighting}[]
\NormalTok{mse }\OtherTok{\textless{}{-}} \FunctionTok{mean}\NormalTok{((predicted }\SpecialCharTok{{-}}\NormalTok{ actual)}\SpecialCharTok{\^{}}\DecValTok{2}\NormalTok{)}
\end{Highlighting}
\end{Shaded}

\subsection{Visualizing Predictions}\label{visualizing-predictions}

\begin{Shaded}
\begin{Highlighting}[]
\FunctionTok{plot}\NormalTok{(}\DecValTok{1}\SpecialCharTok{:}\FunctionTok{length}\NormalTok{(actual), actual, }\AttributeTok{type =} \StringTok{\textquotesingle{}l\textquotesingle{}}\NormalTok{, }\AttributeTok{col =} \StringTok{\textquotesingle{}black\textquotesingle{}}\NormalTok{)}
\FunctionTok{lines}\NormalTok{(}\DecValTok{1}\SpecialCharTok{:}\FunctionTok{length}\NormalTok{(predicted), predicted, }\AttributeTok{col =} \StringTok{\textquotesingle{}red\textquotesingle{}}\NormalTok{)}
\FunctionTok{legend}\NormalTok{(}\StringTok{"topright"}\NormalTok{, }\AttributeTok{legend =} \FunctionTok{c}\NormalTok{(}\StringTok{"Actual"}\NormalTok{, }\StringTok{"Predicted"}\NormalTok{), }\AttributeTok{col =} \FunctionTok{c}\NormalTok{(}\StringTok{"black"}\NormalTok{, }\StringTok{"red"}\NormalTok{), }\AttributeTok{lwd =} \DecValTok{2}\NormalTok{)}
\end{Highlighting}
\end{Shaded}

\section{Example: sin Curve}\label{example-sin-curve}

Generate a sine wave with the function sin(x) for x ranging from 0 to
20, resulting in a total of 400 points. Plot the generated sequence.

\begin{Shaded}
\begin{Highlighting}[]
\FunctionTok{set.seed}\NormalTok{(}\DecValTok{123}\NormalTok{) }\CommentTok{\# set seed}

\NormalTok{data }\OtherTok{\textless{}{-}} \FunctionTok{sin}\NormalTok{(}\FunctionTok{seq}\NormalTok{(}\DecValTok{0}\NormalTok{, }\DecValTok{20}\NormalTok{, }\AttributeTok{length.out =} \DecValTok{400}\NormalTok{)) }\CommentTok{\# create sin data point}

\CommentTok{\# Visualize the sine wave}
\FunctionTok{plot}\NormalTok{(}
  \FunctionTok{as.vector}\NormalTok{(data), }
  \AttributeTok{col =} \StringTok{\textquotesingle{}blue\textquotesingle{}}\NormalTok{, }
  \AttributeTok{type =} \StringTok{\textquotesingle{}l\textquotesingle{}}\NormalTok{, }
  \AttributeTok{ylab =} \StringTok{"Value"}\NormalTok{, }
  \AttributeTok{main =}\StringTok{"Sine Wave"}
\NormalTok{)}
\end{Highlighting}
\end{Shaded}

\includegraphics{RNN-adn-LSTM_files/figure-latex/unnamed-chunk-16-1.pdf}

Divide the cleaned dataset into two parts, the \textbf{last 10 points
for testing y} and \textbf{the rest for training}.

\begin{Shaded}
\begin{Highlighting}[]
\FunctionTok{length}\NormalTok{(data)}
\end{Highlighting}
\end{Shaded}

\begin{verbatim}
## [1] 400
\end{verbatim}

\begin{Shaded}
\begin{Highlighting}[]
\NormalTok{train\_data }\OtherTok{\textless{}{-}}\NormalTok{ data[}\DecValTok{1}\SpecialCharTok{:}\NormalTok{(}\FunctionTok{length}\NormalTok{(data)}\SpecialCharTok{{-}}\DecValTok{10}\NormalTok{)]}
\NormalTok{test\_data }\OtherTok{\textless{}{-}}\NormalTok{ data[(}\FunctionTok{length}\NormalTok{(data)}\SpecialCharTok{{-}}\DecValTok{9}\NormalTok{)}\SpecialCharTok{:}\FunctionTok{length}\NormalTok{(data)]}
\end{Highlighting}
\end{Shaded}

prepares the data in a format compatible with RNNs and LSTMs.

\begin{Shaded}
\begin{Highlighting}[]
\NormalTok{create\_dataset }\OtherTok{\textless{}{-}} \ControlFlowTok{function}\NormalTok{(series, }\AttributeTok{timesteps =} \DecValTok{10}\NormalTok{) \{}
\NormalTok{  x }\OtherTok{\textless{}{-}} \FunctionTok{list}\NormalTok{()}
\NormalTok{  y }\OtherTok{\textless{}{-}} \FunctionTok{c}\NormalTok{()}
  \ControlFlowTok{for}\NormalTok{ (i }\ControlFlowTok{in} \DecValTok{1}\SpecialCharTok{:}\NormalTok{(}\FunctionTok{length}\NormalTok{(series) }\SpecialCharTok{{-}}\NormalTok{ timesteps)) \{}
\NormalTok{    x[[i]] }\OtherTok{\textless{}{-}}\NormalTok{ series[i}\SpecialCharTok{:}\NormalTok{(i }\SpecialCharTok{+}\NormalTok{ timesteps }\SpecialCharTok{{-}} \DecValTok{1}\NormalTok{)]}
\NormalTok{    y }\OtherTok{\textless{}{-}} \FunctionTok{c}\NormalTok{(y, series[i }\SpecialCharTok{+}\NormalTok{ timesteps])}
\NormalTok{  \}}
\NormalTok{  x\_array }\OtherTok{\textless{}{-}} \FunctionTok{do.call}\NormalTok{(rbind, x)}
  \FunctionTok{return}\NormalTok{(}\FunctionTok{list}\NormalTok{(}\AttributeTok{x =}\NormalTok{ x\_array, }\AttributeTok{y =}\NormalTok{ y))}
\NormalTok{\}}
\end{Highlighting}
\end{Shaded}

Create Training Set

\begin{Shaded}
\begin{Highlighting}[]
\NormalTok{timesteps }\OtherTok{\textless{}{-}} \DecValTok{10}
\NormalTok{train\_set }\OtherTok{\textless{}{-}} \FunctionTok{create\_dataset}\NormalTok{(train\_data, timesteps)}
\end{Highlighting}
\end{Shaded}

Convert to Tensor

\begin{Shaded}
\begin{Highlighting}[]
\CommentTok{\#convert to array}
\NormalTok{x\_train\_array }\OtherTok{\textless{}{-}} \FunctionTok{array}\NormalTok{(}\FunctionTok{as.numeric}\NormalTok{(train\_set}\SpecialCharTok{$}\NormalTok{x), }\AttributeTok{dim =} \FunctionTok{c}\NormalTok{(}\FunctionTok{nrow}\NormalTok{(train\_set}\SpecialCharTok{$}\NormalTok{x), timesteps, }\DecValTok{1}\NormalTok{))}
\NormalTok{y\_train\_array }\OtherTok{\textless{}{-}} \FunctionTok{as.numeric}\NormalTok{(train\_set}\SpecialCharTok{$}\NormalTok{y)}

\CommentTok{\# Convert to TensorFlow tensors}
\NormalTok{x\_train }\OtherTok{\textless{}{-}}\NormalTok{ tf}\SpecialCharTok{$}\FunctionTok{convert\_to\_tensor}\NormalTok{(x\_train\_array, }\AttributeTok{dtype =} \StringTok{"float32"}\NormalTok{)}
\NormalTok{y\_train }\OtherTok{\textless{}{-}}\NormalTok{ tf}\SpecialCharTok{$}\FunctionTok{convert\_to\_tensor}\NormalTok{(y\_train\_array, }\AttributeTok{dtype =} \StringTok{"float32"}\NormalTok{)}
\end{Highlighting}
\end{Shaded}

Please first build RNN with length of sequence 10 using the training
data and the RNN method\\
with only one RNN layer with 32 hidden units, and the loss function of
`mse'. You will need to\\
create sequences that the RNN can use to learn the pattern. Each input
sequence should use 10\\
sequential points to predict the next point.

\begin{Shaded}
\begin{Highlighting}[]
\NormalTok{input\_rnn }\OtherTok{\textless{}{-}} \FunctionTok{layer\_input}\NormalTok{(}\AttributeTok{shape =} \FunctionTok{c}\NormalTok{(timesteps, }\DecValTok{1}\NormalTok{))}

\NormalTok{output\_rnn }\OtherTok{\textless{}{-}}\NormalTok{ input\_rnn }\SpecialCharTok{\%\textgreater{}\%}
  \FunctionTok{layer\_simple\_rnn}\NormalTok{(}\AttributeTok{units =} \DecValTok{32}\NormalTok{) }\SpecialCharTok{\%\textgreater{}\%}
  \FunctionTok{layer\_dense}\NormalTok{(}\AttributeTok{units =} \DecValTok{1}\NormalTok{)}

\NormalTok{model\_rnn }\OtherTok{\textless{}{-}} \FunctionTok{keras\_model}\NormalTok{(}\AttributeTok{inputs =}\NormalTok{ input\_rnn, }\AttributeTok{outputs =}\NormalTok{ output\_rnn)}

\NormalTok{model\_rnn}\SpecialCharTok{$}\FunctionTok{compile}\NormalTok{(}
  \AttributeTok{optimizer =} \FunctionTok{optimizer\_adam}\NormalTok{(),}
  \AttributeTok{loss =} \StringTok{"mean\_squared\_error"}
\NormalTok{)}
\end{Highlighting}
\end{Shaded}

Please first build LSTM with length of sequence 10 using the training
data and the LSTM method\\
with only one LSTM layer with 32 hidden units, and the loss function of
`mse'. You will need to\\
create sequences that the RNN can use to learn the pattern. Each input
sequence should use 10\\
sequential points to predict the next point.

\begin{Shaded}
\begin{Highlighting}[]
\NormalTok{input\_lstm }\OtherTok{\textless{}{-}} \FunctionTok{layer\_input}\NormalTok{(}\AttributeTok{shape =} \FunctionTok{c}\NormalTok{(timesteps, }\DecValTok{1}\NormalTok{))}

\NormalTok{output\_lstm }\OtherTok{\textless{}{-}}\NormalTok{ input\_lstm }\SpecialCharTok{\%\textgreater{}\%}
  \FunctionTok{layer\_lstm}\NormalTok{(}\AttributeTok{units =} \DecValTok{32}\NormalTok{) }\SpecialCharTok{\%\textgreater{}\%}
  \FunctionTok{layer\_dense}\NormalTok{(}\AttributeTok{units =} \DecValTok{1}\NormalTok{)}

\NormalTok{model\_lstm }\OtherTok{\textless{}{-}} \FunctionTok{keras\_model}\NormalTok{(}\AttributeTok{inputs =}\NormalTok{ input\_lstm, }\AttributeTok{outputs =}\NormalTok{ output\_lstm)}

\NormalTok{model\_lstm}\SpecialCharTok{$}\FunctionTok{compile}\NormalTok{(}
\AttributeTok{optimizer =} \FunctionTok{optimizer\_adam}\NormalTok{(),}
\AttributeTok{loss =} \StringTok{"mean\_squared\_error"}
\NormalTok{)}
\end{Highlighting}
\end{Shaded}

Plot the training loss for RNN and LSTM.

\begin{Shaded}
\begin{Highlighting}[]
\NormalTok{history\_rnn }\OtherTok{\textless{}{-}}\NormalTok{ model\_rnn}\SpecialCharTok{$}\FunctionTok{fit}\NormalTok{(}
  \AttributeTok{x =}\NormalTok{ x\_train,}
  \AttributeTok{y =}\NormalTok{ y\_train,}
  \AttributeTok{epochs =} \FunctionTok{r\_to\_py}\NormalTok{(}\DecValTok{100}\DataTypeTok{L}\NormalTok{),}
  \AttributeTok{batch\_size =} \FunctionTok{r\_to\_py}\NormalTok{(}\DecValTok{16}\DataTypeTok{L}\NormalTok{),}
  \AttributeTok{verbose =} \DecValTok{0}
\NormalTok{)}

\NormalTok{history\_lstm }\OtherTok{\textless{}{-}}\NormalTok{ model\_lstm}\SpecialCharTok{$}\FunctionTok{fit}\NormalTok{(}
\AttributeTok{x =}\NormalTok{ x\_train,}
\AttributeTok{y =}\NormalTok{ y\_train,}
\AttributeTok{epochs =} \FunctionTok{r\_to\_py}\NormalTok{(}\DecValTok{100}\DataTypeTok{L}\NormalTok{),}
\AttributeTok{batch\_size =} \FunctionTok{r\_to\_py}\NormalTok{(}\DecValTok{16}\DataTypeTok{L}\NormalTok{),}
\AttributeTok{verbose =} \DecValTok{0}
\NormalTok{)}
\end{Highlighting}
\end{Shaded}

\begin{Shaded}
\begin{Highlighting}[]
\FunctionTok{plot}\NormalTok{(}
\NormalTok{  history\_rnn}\SpecialCharTok{$}\NormalTok{history}\SpecialCharTok{$}\NormalTok{loss, }
  \AttributeTok{type =} \StringTok{\textquotesingle{}l\textquotesingle{}}\NormalTok{, }
  \AttributeTok{col =} \StringTok{\textquotesingle{}red\textquotesingle{}}\NormalTok{, }
  \AttributeTok{lwd =} \DecValTok{2}\NormalTok{,}
  \AttributeTok{ylim =} \FunctionTok{range}\NormalTok{(}\FunctionTok{c}\NormalTok{(}
\NormalTok{    history\_rnn}\SpecialCharTok{$}\NormalTok{history}\SpecialCharTok{$}\NormalTok{loss,}
\NormalTok{    history\_lstm}\SpecialCharTok{$}\NormalTok{history}\SpecialCharTok{$}\NormalTok{loss)}
\NormalTok{    ),}
  \AttributeTok{xlab =} \StringTok{"Epoch"}\NormalTok{, }
  \AttributeTok{ylab =} \StringTok{"Training Loss"}\NormalTok{, }
  \AttributeTok{main =} \StringTok{"Training Loss Comparison"}
\NormalTok{)}

\FunctionTok{lines}\NormalTok{(}
\NormalTok{  history\_lstm}\SpecialCharTok{$}\NormalTok{history}\SpecialCharTok{$}\NormalTok{loss, }
  \AttributeTok{col =} \StringTok{\textquotesingle{}blue\textquotesingle{}}\NormalTok{, }
  \AttributeTok{lwd =} \DecValTok{2}
\NormalTok{)}

\FunctionTok{legend}\NormalTok{(}
  \StringTok{"topright"}\NormalTok{, }
  \AttributeTok{legend =} \FunctionTok{c}\NormalTok{(}\StringTok{"RNN"}\NormalTok{, }\StringTok{"LSTM"}\NormalTok{), }
  \AttributeTok{col =} \FunctionTok{c}\NormalTok{(}\StringTok{"red"}\NormalTok{, }\StringTok{"blue"}\NormalTok{),}
  \AttributeTok{lwd =} \DecValTok{2}
\NormalTok{)}
\end{Highlighting}
\end{Shaded}

\includegraphics{RNN-adn-LSTM_files/figure-latex/unnamed-chunk-24-1.pdf}

Please compute the MSE of each model in Question 3 and Question 4 for
the last 10 points.\\
Compare the performance of the algorithms in Question 3 and Question 4
in length-10 forecast\\
using the testing data, and plot the predictions along with the actual
data.

\begin{Shaded}
\begin{Highlighting}[]
\CommentTok{\# Predictions}
\NormalTok{input\_seq }\OtherTok{\textless{}{-}}\NormalTok{ data[}\DecValTok{381}\SpecialCharTok{:}\DecValTok{390}\NormalTok{]}

\NormalTok{input\_seq\_array }\OtherTok{\textless{}{-}} \FunctionTok{array}\NormalTok{(}
  \FunctionTok{as.numeric}\NormalTok{(input\_seq), }
  \AttributeTok{dim =} \FunctionTok{c}\NormalTok{(}\DecValTok{1}\NormalTok{, timesteps,}\DecValTok{1}\NormalTok{)}
\NormalTok{)}

\NormalTok{input\_seq\_array }\OtherTok{\textless{}{-}}\NormalTok{ tf}\SpecialCharTok{$}\FunctionTok{convert\_to\_tensor}\NormalTok{(}
\NormalTok{  input\_seq\_array, }
  \AttributeTok{dtype =} \StringTok{"float32"}
\NormalTok{)}

\NormalTok{predict\_next }\OtherTok{\textless{}{-}} \ControlFlowTok{function}\NormalTok{(model, input\_seq\_array, }\AttributeTok{n\_steps =} \DecValTok{10}\NormalTok{) \{}
\NormalTok{  preds }\OtherTok{\textless{}{-}} \FunctionTok{c}\NormalTok{()}
\NormalTok{  current\_input }\OtherTok{\textless{}{-}}\NormalTok{ input\_seq\_array}
  
  \ControlFlowTok{for}\NormalTok{ (i }\ControlFlowTok{in} \DecValTok{1}\SpecialCharTok{:}\NormalTok{n\_steps) \{}
\NormalTok{    pred }\OtherTok{\textless{}{-}}\NormalTok{ model}\SpecialCharTok{$}\FunctionTok{predict}\NormalTok{(current\_input)}
\NormalTok{    preds }\OtherTok{\textless{}{-}} \FunctionTok{c}\NormalTok{(preds, }\FunctionTok{as.numeric}\NormalTok{(pred))}
\NormalTok{    updated\_seq }\OtherTok{\textless{}{-}} \FunctionTok{c}\NormalTok{(}\FunctionTok{as.numeric}\NormalTok{(current\_input)[}\SpecialCharTok{{-}}\DecValTok{1}\NormalTok{], }\FunctionTok{as.numeric}\NormalTok{(pred))}
\NormalTok{    current\_input }\OtherTok{\textless{}{-}} \FunctionTok{array}\NormalTok{(updated\_seq, }\AttributeTok{dim =} \FunctionTok{c}\NormalTok{(}\DecValTok{1}\NormalTok{, timesteps, }\DecValTok{1}\NormalTok{))}
\NormalTok{    current\_input }\OtherTok{\textless{}{-}}\NormalTok{ tf}\SpecialCharTok{$}\FunctionTok{convert\_to\_tensor}\NormalTok{(current\_input, }\AttributeTok{dtype =} \StringTok{"float32"}\NormalTok{)}
\NormalTok{  \}}
  
  \FunctionTok{return}\NormalTok{(preds)}
\NormalTok{\}}
\end{Highlighting}
\end{Shaded}

\begin{Shaded}
\begin{Highlighting}[]
\CommentTok{\# 10{-}step prediction}
\NormalTok{rnn\_preds }\OtherTok{\textless{}{-}} \FunctionTok{predict\_next}\NormalTok{(}
\NormalTok{  model\_rnn, }
\NormalTok{  input\_seq\_array, }
  \AttributeTok{n\_steps =} \DecValTok{10}
\NormalTok{)}
\end{Highlighting}
\end{Shaded}

\begin{verbatim}
## 1/1 ━━━━━━━━━━━━━━━━━━━━ 0s 43ms/step1/1 ━━━━━━━━━━━━━━━━━━━━ 0s 49ms/step
## 1/1 ━━━━━━━━━━━━━━━━━━━━ 0s 6ms/step1/1 ━━━━━━━━━━━━━━━━━━━━ 0s 12ms/step
## 1/1 ━━━━━━━━━━━━━━━━━━━━ 0s 6ms/step1/1 ━━━━━━━━━━━━━━━━━━━━ 0s 12ms/step
## 1/1 ━━━━━━━━━━━━━━━━━━━━ 0s 6ms/step1/1 ━━━━━━━━━━━━━━━━━━━━ 0s 13ms/step
## 1/1 ━━━━━━━━━━━━━━━━━━━━ 0s 6ms/step1/1 ━━━━━━━━━━━━━━━━━━━━ 0s 12ms/step
## 1/1 ━━━━━━━━━━━━━━━━━━━━ 0s 6ms/step1/1 ━━━━━━━━━━━━━━━━━━━━ 0s 12ms/step
## 1/1 ━━━━━━━━━━━━━━━━━━━━ 0s 6ms/step1/1 ━━━━━━━━━━━━━━━━━━━━ 0s 12ms/step
## 1/1 ━━━━━━━━━━━━━━━━━━━━ 0s 7ms/step1/1 ━━━━━━━━━━━━━━━━━━━━ 0s 12ms/step
## 1/1 ━━━━━━━━━━━━━━━━━━━━ 0s 6ms/step1/1 ━━━━━━━━━━━━━━━━━━━━ 0s 12ms/step
## 1/1 ━━━━━━━━━━━━━━━━━━━━ 0s 6ms/step1/1 ━━━━━━━━━━━━━━━━━━━━ 0s 12ms/step
\end{verbatim}

\begin{Shaded}
\begin{Highlighting}[]
\NormalTok{lstm\_preds }\OtherTok{\textless{}{-}} \FunctionTok{predict\_next}\NormalTok{(}
\NormalTok{  model\_lstm, }
\NormalTok{  input\_seq\_array, }
  \AttributeTok{n\_steps =} \DecValTok{10}
\NormalTok{)}
\end{Highlighting}
\end{Shaded}

\begin{verbatim}
## 1/1 ━━━━━━━━━━━━━━━━━━━━ 0s 48ms/step1/1 ━━━━━━━━━━━━━━━━━━━━ 0s 54ms/step
## 1/1 ━━━━━━━━━━━━━━━━━━━━ 0s 6ms/step1/1 ━━━━━━━━━━━━━━━━━━━━ 0s 12ms/step
## 1/1 ━━━━━━━━━━━━━━━━━━━━ 0s 7ms/step1/1 ━━━━━━━━━━━━━━━━━━━━ 0s 12ms/step
## 1/1 ━━━━━━━━━━━━━━━━━━━━ 0s 6ms/step1/1 ━━━━━━━━━━━━━━━━━━━━ 0s 12ms/step
## 1/1 ━━━━━━━━━━━━━━━━━━━━ 0s 6ms/step1/1 ━━━━━━━━━━━━━━━━━━━━ 0s 12ms/step
## 1/1 ━━━━━━━━━━━━━━━━━━━━ 0s 7ms/step1/1 ━━━━━━━━━━━━━━━━━━━━ 0s 12ms/step
## 1/1 ━━━━━━━━━━━━━━━━━━━━ 0s 6ms/step1/1 ━━━━━━━━━━━━━━━━━━━━ 0s 12ms/step
## 1/1 ━━━━━━━━━━━━━━━━━━━━ 0s 6ms/step1/1 ━━━━━━━━━━━━━━━━━━━━ 0s 12ms/step
## 1/1 ━━━━━━━━━━━━━━━━━━━━ 0s 6ms/step1/1 ━━━━━━━━━━━━━━━━━━━━ 0s 12ms/step
## 1/1 ━━━━━━━━━━━━━━━━━━━━ 0s 6ms/step1/1 ━━━━━━━━━━━━━━━━━━━━ 0s 12ms/step
\end{verbatim}

\subsubsection{Comparison}\label{comparison}

\begin{Shaded}
\begin{Highlighting}[]
\FunctionTok{plot}\NormalTok{(}
  \DecValTok{1}\SpecialCharTok{:}\DecValTok{400}\NormalTok{, }
\NormalTok{  data, }
  \AttributeTok{type =} \StringTok{\textquotesingle{}l\textquotesingle{}}\NormalTok{, }
  \AttributeTok{col =} \StringTok{\textquotesingle{}black\textquotesingle{}}\NormalTok{, }
  \AttributeTok{lwd =} \DecValTok{2}\NormalTok{, }
  \AttributeTok{ylab =} \StringTok{"Value"}\NormalTok{,}
  \AttributeTok{xlab =} \StringTok{"Time Step"}\NormalTok{,}
  \AttributeTok{main =} \StringTok{"True vs Predicted (Last 10 Points)"}
\NormalTok{)}

\FunctionTok{lines}\NormalTok{(}
  \DecValTok{391}\SpecialCharTok{:}\DecValTok{400}\NormalTok{, }
\NormalTok{  rnn\_preds, }
  \AttributeTok{col =} \StringTok{\textquotesingle{}red\textquotesingle{}}\NormalTok{, }
  \AttributeTok{lwd =} \DecValTok{2}
\NormalTok{)}

\FunctionTok{lines}\NormalTok{(}
  \DecValTok{391}\SpecialCharTok{:}\DecValTok{400}\NormalTok{, }
\NormalTok{  lstm\_preds, }
  \AttributeTok{col =} \StringTok{\textquotesingle{}blue\textquotesingle{}}\NormalTok{, }
  \AttributeTok{lwd =} \DecValTok{2}
\NormalTok{)}

\FunctionTok{points}\NormalTok{(}
  \DecValTok{391}\SpecialCharTok{:}\DecValTok{400}\NormalTok{, }
\NormalTok{  test\_data, }
  \AttributeTok{col =} \StringTok{\textquotesingle{}black\textquotesingle{}}\NormalTok{, }
  \AttributeTok{pch =} \DecValTok{16}
\NormalTok{)}

\FunctionTok{legend}\NormalTok{(}
  \StringTok{"bottomleft"}\NormalTok{, }
  \AttributeTok{legend =} \FunctionTok{c}\NormalTok{(}\StringTok{"True"}\NormalTok{, }\StringTok{"RNN Prediction"}\NormalTok{, }\StringTok{"LSTM Prediction"}\NormalTok{),}
  \AttributeTok{col =} \FunctionTok{c}\NormalTok{(}\StringTok{"black"}\NormalTok{, }\StringTok{"red"}\NormalTok{, }\StringTok{"blue"}\NormalTok{), }
  \AttributeTok{lwd =} \DecValTok{2}\NormalTok{, }
  \AttributeTok{pch =} \FunctionTok{c}\NormalTok{(}\DecValTok{16}\NormalTok{, }\ConstantTok{NA}\NormalTok{, }\ConstantTok{NA}\NormalTok{)}
\NormalTok{)}
\end{Highlighting}
\end{Shaded}

\includegraphics{RNN-adn-LSTM_files/figure-latex/unnamed-chunk-28-1.pdf}

\section{Example: AAPL}\label{example-aapl}

\begin{Shaded}
\begin{Highlighting}[]
\NormalTok{data }\OtherTok{\textless{}{-}} \FunctionTok{read.csv}\NormalTok{(}\StringTok{"AAPL.csv"}\NormalTok{)}

\CommentTok{\# transform the data type from \textquotesingle{}chr\textquotesingle{} to \textquotesingle{}Date\textquotesingle{}}
\NormalTok{data}\SpecialCharTok{$}\NormalTok{Date }\OtherTok{=} \FunctionTok{as.Date}\NormalTok{(data}\SpecialCharTok{$}\NormalTok{Date)}

\FunctionTok{head}\NormalTok{(data)}
\end{Highlighting}
\end{Shaded}

\begin{verbatim}
##         Date   Open   High    Low  Close Adj.Close    Volume
## 1 2021-03-29 121.65 122.58 120.73 121.39  120.6728  80819200
## 2 2021-03-30 120.11 120.40 118.86 119.90  119.1916  85671900
## 3 2021-03-31 121.65 123.52 121.15 122.15  121.4283 118323800
## 4 2021-04-01 123.66 124.18 122.49 123.00  122.2733  75089100
## 5 2021-04-05 123.87 126.16 123.07 125.90  125.1561  88651200
## 6 2021-04-06 126.50 127.13 125.65 126.21  125.4643  80171300
\end{verbatim}

Plot the close price vs.~date to visualize the data we will analyze.

\begin{Shaded}
\begin{Highlighting}[]
\FunctionTok{ggplot}\NormalTok{(}
\NormalTok{  data, }
  \FunctionTok{aes}\NormalTok{(}\AttributeTok{x=}\NormalTok{Date, }\AttributeTok{y =}\NormalTok{ Close)}
\NormalTok{) }\SpecialCharTok{+} \FunctionTok{geom\_line}\NormalTok{()}
\end{Highlighting}
\end{Shaded}

\includegraphics{RNN-adn-LSTM_files/figure-latex/unnamed-chunk-30-1.pdf}

Then, use the `min-max scaler' to normalize our stock price data.

Scaled\_x = (x-lagged\_min(x))/(lagged\_max(x) -- lagged\_min(x)).

Please report the values of the last 10 scaled close prices.

\begin{Shaded}
\begin{Highlighting}[]
\NormalTok{data}\SpecialCharTok{$}\NormalTok{min\_lagged }\OtherTok{=} \FunctionTok{lag}\NormalTok{(data}\SpecialCharTok{$}\NormalTok{Low)}
\NormalTok{data}\SpecialCharTok{$}\NormalTok{max\_lagged }\OtherTok{=} \FunctionTok{lag}\NormalTok{(data}\SpecialCharTok{$}\NormalTok{High)}

\NormalTok{data}\SpecialCharTok{$}\NormalTok{Close\_norm }\OtherTok{=}\NormalTok{ (data}\SpecialCharTok{$}\NormalTok{Close }\SpecialCharTok{{-}}\NormalTok{ data}\SpecialCharTok{$}\NormalTok{min\_lagged) }\SpecialCharTok{/}\NormalTok{ (data}\SpecialCharTok{$}\NormalTok{max\_lagged }\SpecialCharTok{{-}}\NormalTok{data}\SpecialCharTok{$}\NormalTok{min\_lagged)}
\end{Highlighting}
\end{Shaded}

\begin{Shaded}
\begin{Highlighting}[]
\NormalTok{model\_data }\OtherTok{\textless{}{-}} \FunctionTok{matrix}\NormalTok{(}
\NormalTok{  data}\SpecialCharTok{$}\NormalTok{Close\_norm[}\SpecialCharTok{{-}}\DecValTok{1}\NormalTok{] }\CommentTok{\#excluding volume}
\NormalTok{)}
\end{Highlighting}
\end{Shaded}

Divide the cleaned dataset into two parts, the last 10 prices for
testing y and the rest for\\
training (the last 10 prices in the training set will be used as testing
x). Please report how many\\
days of stock price are divided into the training set.

\begin{Shaded}
\begin{Highlighting}[]
\FunctionTok{set.seed}\NormalTok{(}\DecValTok{123}\NormalTok{)}

\CommentTok{\# Create sliding window dataset}
\NormalTok{create\_dataset }\OtherTok{\textless{}{-}} \ControlFlowTok{function}\NormalTok{(series, }\AttributeTok{timesteps =} \DecValTok{10}\NormalTok{) \{}
\NormalTok{  x }\OtherTok{\textless{}{-}} \FunctionTok{list}\NormalTok{()}
\NormalTok{  y }\OtherTok{\textless{}{-}} \FunctionTok{c}\NormalTok{()}
  \ControlFlowTok{for}\NormalTok{ (i }\ControlFlowTok{in} \DecValTok{1}\SpecialCharTok{:}\NormalTok{(}\FunctionTok{length}\NormalTok{(series) }\SpecialCharTok{{-}}\NormalTok{ timesteps)) \{}
\NormalTok{  x[[i]] }\OtherTok{\textless{}{-}}\NormalTok{ series[i}\SpecialCharTok{:}\NormalTok{(i }\SpecialCharTok{+}\NormalTok{ timesteps }\SpecialCharTok{{-}} \DecValTok{1}\NormalTok{)]}
\NormalTok{  y }\OtherTok{\textless{}{-}} \FunctionTok{c}\NormalTok{(y, series[i }\SpecialCharTok{+}\NormalTok{ timesteps])}
\NormalTok{  \}}
\NormalTok{  x\_array }\OtherTok{\textless{}{-}} \FunctionTok{do.call}\NormalTok{(rbind, x)}
  \FunctionTok{return}\NormalTok{(}\FunctionTok{list}\NormalTok{(}\AttributeTok{x =}\NormalTok{ x\_array, }\AttributeTok{y =}\NormalTok{ y))}
\NormalTok{\}}

\CommentTok{\# Prepare training data}
\NormalTok{timesteps }\OtherTok{\textless{}{-}} \DecValTok{10}

\NormalTok{train\_data }\OtherTok{\textless{}{-}}\NormalTok{ model\_data[}\DecValTok{1}\SpecialCharTok{:}\NormalTok{(}\FunctionTok{length}\NormalTok{(model\_data)}\SpecialCharTok{{-}}\DecValTok{10}\NormalTok{)] }\CommentTok{\# training = all except last 10 days}

\NormalTok{test\_data }\OtherTok{\textless{}{-}}\NormalTok{ model\_data[(}\FunctionTok{length}\NormalTok{(model\_data)}\SpecialCharTok{{-}}\DecValTok{9}\NormalTok{)}\SpecialCharTok{:}\FunctionTok{length}\NormalTok{(model\_data)] }\CommentTok{\# true last 10 points}

\NormalTok{train\_set }\OtherTok{\textless{}{-}} \FunctionTok{create\_dataset}\NormalTok{(train\_data, timesteps)}

\NormalTok{x\_train\_array }\OtherTok{\textless{}{-}} \FunctionTok{array}\NormalTok{(}
  \FunctionTok{as.numeric}\NormalTok{(train\_set}\SpecialCharTok{$}\NormalTok{x), }
  \AttributeTok{dim =} \FunctionTok{c}\NormalTok{(}\FunctionTok{nrow}\NormalTok{(train\_set}\SpecialCharTok{$}\NormalTok{x), timesteps, }\DecValTok{1}\NormalTok{)}
\NormalTok{)}

\NormalTok{y\_train\_array }\OtherTok{\textless{}{-}} \FunctionTok{as.numeric}\NormalTok{(train\_set}\SpecialCharTok{$}\NormalTok{y)}

\NormalTok{x\_train }\OtherTok{\textless{}{-}}\NormalTok{ tf}\SpecialCharTok{$}\FunctionTok{convert\_to\_tensor}\NormalTok{(}
\NormalTok{  x\_train\_array, }
  \AttributeTok{dtype =} \StringTok{"float32"}
\NormalTok{)}

\NormalTok{y\_train }\OtherTok{\textless{}{-}}\NormalTok{ tf}\SpecialCharTok{$}\FunctionTok{convert\_to\_tensor}\NormalTok{(}
\NormalTok{  y\_train\_array, }
  \AttributeTok{dtype =} \StringTok{"float32"}
\NormalTok{)}
\end{Highlighting}
\end{Shaded}

Please first build the predictive model to predict 10-days stock close
price using the training\\
data and the LSTM method with only one LSTM layer with 200 hidden units,
and the loss\\
function of `mse'. Please make predictions on the 10 observations in the
testing set by using the\\
last 10 in the training dataset and compute the Test MSE using the
testing data. Scale the\\
predicted stock price back and plot the 10-days predictions and the true
stock close price in the\\
same figure. Also, try to predict the close price of 2022-3-29.

\begin{Shaded}
\begin{Highlighting}[]
\CommentTok{\# Build RNN Model}
\NormalTok{input\_rnn }\OtherTok{\textless{}{-}} \FunctionTok{layer\_input}\NormalTok{(}\AttributeTok{shape =} \FunctionTok{c}\NormalTok{(timesteps, }\DecValTok{1}\NormalTok{))}

\NormalTok{output\_rnn }\OtherTok{\textless{}{-}}\NormalTok{ input\_rnn }\SpecialCharTok{\%\textgreater{}\%}
  \FunctionTok{layer\_simple\_rnn}\NormalTok{(}\AttributeTok{units =} \DecValTok{32}\NormalTok{) }\SpecialCharTok{\%\textgreater{}\%}
  \FunctionTok{layer\_dense}\NormalTok{(}\AttributeTok{units =} \DecValTok{1}\NormalTok{)}

\NormalTok{model\_rnn }\OtherTok{\textless{}{-}} \FunctionTok{keras\_model}\NormalTok{(}\AttributeTok{inputs =}\NormalTok{ input\_rnn, }\AttributeTok{outputs =}\NormalTok{ output\_rnn)}

\NormalTok{model\_rnn}\SpecialCharTok{$}\FunctionTok{compile}\NormalTok{(}
  \AttributeTok{optimizer =} \FunctionTok{optimizer\_adam}\NormalTok{(),}
  \AttributeTok{loss =} \StringTok{"mean\_squared\_error"}
\NormalTok{)}

\NormalTok{history\_rnn }\OtherTok{\textless{}{-}}\NormalTok{ model\_rnn}\SpecialCharTok{$}\FunctionTok{fit}\NormalTok{(}
  \AttributeTok{x =}\NormalTok{ x\_train,}
  \AttributeTok{y =}\NormalTok{ y\_train,}
  \AttributeTok{epochs =} \FunctionTok{r\_to\_py}\NormalTok{(}\DecValTok{100}\DataTypeTok{L}\NormalTok{),}
  \AttributeTok{batch\_size =} \FunctionTok{r\_to\_py}\NormalTok{(}\DecValTok{16}\DataTypeTok{L}\NormalTok{),}
  \AttributeTok{verbose =} \DecValTok{0}
\NormalTok{)}
\end{Highlighting}
\end{Shaded}

Please first build the predictive model to predict 10-days stock close
price using the training\\
data and the RNN method with only one RNN layer with 200 hidden units,
and the loss function\\
of `mse'. Please make predictions on the 10 observations in the testing
set by using the last 10\\
in the training dataset and compute the Test MSE using the testing data.
Scale the predicted\\
stock price back and plot the 10-days predictions and the true stock
close price in the same\\
figure. Also, try to predict the close price of 2022-3-29.

\begin{Shaded}
\begin{Highlighting}[]
\CommentTok{\# Build LSTM Model}
\NormalTok{input\_lstm }\OtherTok{\textless{}{-}} \FunctionTok{layer\_input}\NormalTok{(}\AttributeTok{shape =} \FunctionTok{c}\NormalTok{(timesteps, }\DecValTok{1}\NormalTok{))}

\NormalTok{output\_lstm }\OtherTok{\textless{}{-}}\NormalTok{ input\_lstm }\SpecialCharTok{\%\textgreater{}\%}
  \FunctionTok{layer\_lstm}\NormalTok{(}\AttributeTok{units =} \DecValTok{32}\NormalTok{) }\SpecialCharTok{\%\textgreater{}\%}
  \FunctionTok{layer\_dense}\NormalTok{(}\AttributeTok{units =} \DecValTok{1}\NormalTok{)}

\NormalTok{model\_lstm }\OtherTok{\textless{}{-}} \FunctionTok{keras\_model}\NormalTok{(}\AttributeTok{inputs =}\NormalTok{ input\_lstm, }\AttributeTok{outputs =}\NormalTok{ output\_lstm)}

\NormalTok{model\_lstm}\SpecialCharTok{$}\FunctionTok{compile}\NormalTok{(}
  \AttributeTok{optimizer =} \FunctionTok{optimizer\_adam}\NormalTok{(),}
  \AttributeTok{loss =} \StringTok{"mean\_squared\_error"}
\NormalTok{)}

\NormalTok{history\_lstm }\OtherTok{\textless{}{-}}\NormalTok{ model\_lstm}\SpecialCharTok{$}\FunctionTok{fit}\NormalTok{(}
  \AttributeTok{x =}\NormalTok{ x\_train,}
  \AttributeTok{y =}\NormalTok{ y\_train,}
  \AttributeTok{epochs =} \FunctionTok{r\_to\_py}\NormalTok{(}\DecValTok{100}\DataTypeTok{L}\NormalTok{),}
  \AttributeTok{batch\_size =} \FunctionTok{r\_to\_py}\NormalTok{(}\DecValTok{16}\DataTypeTok{L}\NormalTok{),}
  \AttributeTok{verbose =} \DecValTok{0}
\NormalTok{)}
\end{Highlighting}
\end{Shaded}

\subsubsection{Comparison}\label{comparison-1}

Please compare the performance of the algorithms in Question 4 and
Question 5 in 10-days\\
forecast using the testing data

\begin{Shaded}
\begin{Highlighting}[]
\CommentTok{\# Prepare input for forecasting}
\NormalTok{input\_seq }\OtherTok{\textless{}{-}}\NormalTok{ model\_data[(}\FunctionTok{length}\NormalTok{(model\_data)}\SpecialCharTok{{-}}\DecValTok{10}\NormalTok{)}\SpecialCharTok{:}\NormalTok{(}\FunctionTok{length}\NormalTok{(model\_data)}\SpecialCharTok{{-}}\DecValTok{1}\NormalTok{)] }\CommentTok{\# 10 points before test}

\NormalTok{input\_seq\_array }\OtherTok{\textless{}{-}} \FunctionTok{array}\NormalTok{(}\FunctionTok{as.numeric}\NormalTok{(input\_seq), }\AttributeTok{dim =} \FunctionTok{c}\NormalTok{(}\DecValTok{1}\NormalTok{, timesteps,}\DecValTok{1}\NormalTok{))}

\NormalTok{input\_seq\_array }\OtherTok{\textless{}{-}}\NormalTok{ tf}\SpecialCharTok{$}\FunctionTok{convert\_to\_tensor}\NormalTok{(input\_seq\_array, }\AttributeTok{dtype =} \StringTok{"float32"}\NormalTok{)}

\CommentTok{\# Prediction function}
\NormalTok{predict\_next }\OtherTok{\textless{}{-}} \ControlFlowTok{function}\NormalTok{(model, input\_seq\_array, }\AttributeTok{n\_steps =} \DecValTok{10}\NormalTok{) \{}
\NormalTok{  preds }\OtherTok{\textless{}{-}} \FunctionTok{c}\NormalTok{()}
\NormalTok{  current\_input }\OtherTok{\textless{}{-}}\NormalTok{ input\_seq\_array}
  \ControlFlowTok{for}\NormalTok{ (i }\ControlFlowTok{in} \DecValTok{1}\SpecialCharTok{:}\NormalTok{n\_steps) \{}
\NormalTok{    pred }\OtherTok{\textless{}{-}}\NormalTok{ model}\SpecialCharTok{$}\FunctionTok{predict}\NormalTok{(current\_input)}
\NormalTok{    preds }\OtherTok{\textless{}{-}} \FunctionTok{c}\NormalTok{(preds, }\FunctionTok{as.numeric}\NormalTok{(pred))}
\NormalTok{    updated\_seq }\OtherTok{\textless{}{-}} \FunctionTok{c}\NormalTok{(}\FunctionTok{as.numeric}\NormalTok{(current\_input)[}\SpecialCharTok{{-}}\DecValTok{1}\NormalTok{], }\FunctionTok{as.numeric}\NormalTok{(pred))}
\NormalTok{    current\_input }\OtherTok{\textless{}{-}} \FunctionTok{array}\NormalTok{(updated\_seq, }\AttributeTok{dim =} \FunctionTok{c}\NormalTok{(}\DecValTok{1}\NormalTok{, timesteps, }\DecValTok{1}\NormalTok{))}
\NormalTok{    current\_input }\OtherTok{\textless{}{-}}\NormalTok{ tf}\SpecialCharTok{$}\FunctionTok{convert\_to\_tensor}\NormalTok{(current\_input, }\AttributeTok{dtype =} \StringTok{"float32"}\NormalTok{)}
\NormalTok{  \}}
  \FunctionTok{return}\NormalTok{(preds)}
\NormalTok{\}}
\end{Highlighting}
\end{Shaded}

\begin{Shaded}
\begin{Highlighting}[]
\NormalTok{rnn\_preds\_norm }\OtherTok{\textless{}{-}} \FunctionTok{predict\_next}\NormalTok{(}
\NormalTok{  model\_rnn, }
\NormalTok{  input\_seq\_array, }
  \AttributeTok{n\_steps =} \DecValTok{10}
\NormalTok{)}
\end{Highlighting}
\end{Shaded}

\begin{verbatim}
## 1/1 ━━━━━━━━━━━━━━━━━━━━ 0s 42ms/step1/1 ━━━━━━━━━━━━━━━━━━━━ 0s 48ms/step
## 1/1 ━━━━━━━━━━━━━━━━━━━━ 0s 6ms/step1/1 ━━━━━━━━━━━━━━━━━━━━ 0s 12ms/step
## 1/1 ━━━━━━━━━━━━━━━━━━━━ 0s 6ms/step1/1 ━━━━━━━━━━━━━━━━━━━━ 0s 12ms/step
## 1/1 ━━━━━━━━━━━━━━━━━━━━ 0s 6ms/step1/1 ━━━━━━━━━━━━━━━━━━━━ 0s 12ms/step
## 1/1 ━━━━━━━━━━━━━━━━━━━━ 0s 6ms/step1/1 ━━━━━━━━━━━━━━━━━━━━ 0s 12ms/step
## 1/1 ━━━━━━━━━━━━━━━━━━━━ 0s 6ms/step1/1 ━━━━━━━━━━━━━━━━━━━━ 0s 12ms/step
## 1/1 ━━━━━━━━━━━━━━━━━━━━ 0s 6ms/step1/1 ━━━━━━━━━━━━━━━━━━━━ 0s 12ms/step
## 1/1 ━━━━━━━━━━━━━━━━━━━━ 0s 7ms/step1/1 ━━━━━━━━━━━━━━━━━━━━ 0s 12ms/step
## 1/1 ━━━━━━━━━━━━━━━━━━━━ 0s 6ms/step1/1 ━━━━━━━━━━━━━━━━━━━━ 0s 12ms/step
## 1/1 ━━━━━━━━━━━━━━━━━━━━ 0s 6ms/step1/1 ━━━━━━━━━━━━━━━━━━━━ 0s 12ms/step
\end{verbatim}

\begin{Shaded}
\begin{Highlighting}[]
\NormalTok{lstm\_preds\_norm }\OtherTok{\textless{}{-}} \FunctionTok{predict\_next}\NormalTok{(}
\NormalTok{  model\_lstm, }
\NormalTok{  input\_seq\_array, }
  \AttributeTok{n\_steps =} \DecValTok{10}
\NormalTok{)}
\end{Highlighting}
\end{Shaded}

\begin{verbatim}
## 1/1 ━━━━━━━━━━━━━━━━━━━━ 0s 48ms/step1/1 ━━━━━━━━━━━━━━━━━━━━ 0s 53ms/step
## 1/1 ━━━━━━━━━━━━━━━━━━━━ 0s 7ms/step1/1 ━━━━━━━━━━━━━━━━━━━━ 0s 12ms/step
## 1/1 ━━━━━━━━━━━━━━━━━━━━ 0s 6ms/step1/1 ━━━━━━━━━━━━━━━━━━━━ 0s 12ms/step
## 1/1 ━━━━━━━━━━━━━━━━━━━━ 0s 6ms/step1/1 ━━━━━━━━━━━━━━━━━━━━ 0s 12ms/step
## 1/1 ━━━━━━━━━━━━━━━━━━━━ 0s 6ms/step1/1 ━━━━━━━━━━━━━━━━━━━━ 0s 12ms/step
## 1/1 ━━━━━━━━━━━━━━━━━━━━ 0s 6ms/step1/1 ━━━━━━━━━━━━━━━━━━━━ 0s 12ms/step
## 1/1 ━━━━━━━━━━━━━━━━━━━━ 0s 6ms/step1/1 ━━━━━━━━━━━━━━━━━━━━ 0s 12ms/step
## 1/1 ━━━━━━━━━━━━━━━━━━━━ 0s 7ms/step1/1 ━━━━━━━━━━━━━━━━━━━━ 0s 12ms/step
## 1/1 ━━━━━━━━━━━━━━━━━━━━ 0s 6ms/step1/1 ━━━━━━━━━━━━━━━━━━━━ 0s 12ms/step
## 1/1 ━━━━━━━━━━━━━━━━━━━━ 0s 6ms/step1/1 ━━━━━━━━━━━━━━━━━━━━ 0s 12ms/step
\end{verbatim}

Scale Back

\begin{Shaded}
\begin{Highlighting}[]
\CommentTok{\# Initialize vectors to store scaled{-}back predictions and true prices}
\NormalTok{rnn\_preds\_original }\OtherTok{\textless{}{-}} \FunctionTok{numeric}\NormalTok{(}\DecValTok{10}\NormalTok{)}
\NormalTok{lstm\_preds\_original }\OtherTok{\textless{}{-}} \FunctionTok{numeric}\NormalTok{(}\DecValTok{10}\NormalTok{)}
\NormalTok{true\_prices\_original }\OtherTok{\textless{}{-}} \FunctionTok{numeric}\NormalTok{(}\DecValTok{10}\NormalTok{)}

\CommentTok{\# Scale back one by one (because each day uses its own lagged High/Low)}
\ControlFlowTok{for}\NormalTok{ (i }\ControlFlowTok{in} \DecValTok{1}\SpecialCharTok{:}\DecValTok{10}\NormalTok{) \{}
\NormalTok{  lagged\_min }\OtherTok{\textless{}{-}}\NormalTok{ data}\SpecialCharTok{$}\NormalTok{Low[}\FunctionTok{length}\NormalTok{(train\_data) }\SpecialCharTok{+}\NormalTok{ i }\SpecialCharTok{{-}} \DecValTok{1}\NormalTok{] }\CommentTok{\# lag(Low) for that day}
\NormalTok{  lagged\_max }\OtherTok{\textless{}{-}}\NormalTok{ data}\SpecialCharTok{$}\NormalTok{High[}\FunctionTok{length}\NormalTok{(train\_data) }\SpecialCharTok{+}\NormalTok{ i }\SpecialCharTok{{-}} \DecValTok{1}\NormalTok{] }\CommentTok{\# lag(High) for that day}
\NormalTok{  rnn\_preds\_original[i] }\OtherTok{\textless{}{-}}\NormalTok{ rnn\_preds\_norm[i] }\SpecialCharTok{*}\NormalTok{ (lagged\_max }\SpecialCharTok{{-}}\NormalTok{ lagged\_min) }\SpecialCharTok{+}\NormalTok{ lagged\_min}
\NormalTok{  lstm\_preds\_original[i] }\OtherTok{\textless{}{-}}\NormalTok{ lstm\_preds\_norm[i] }\SpecialCharTok{*}\NormalTok{ (lagged\_max }\SpecialCharTok{{-}}\NormalTok{ lagged\_min) }\SpecialCharTok{+}\NormalTok{ lagged\_min}
\NormalTok{  true\_prices\_original[i] }\OtherTok{\textless{}{-}}\NormalTok{ data}\SpecialCharTok{$}\NormalTok{Close[}\FunctionTok{length}\NormalTok{(train\_data) }\SpecialCharTok{+}\NormalTok{ i] }\CommentTok{\# true Close at day}
\NormalTok{\}}
\end{Highlighting}
\end{Shaded}

MSE

\begin{Shaded}
\begin{Highlighting}[]
\NormalTok{mse\_rnn }\OtherTok{\textless{}{-}} \FunctionTok{mean}\NormalTok{((rnn\_preds\_original }\SpecialCharTok{{-}}\NormalTok{ true\_prices\_original)}\SpecialCharTok{\^{}}\DecValTok{2}\NormalTok{)}

\NormalTok{mse\_lstm }\OtherTok{\textless{}{-}} \FunctionTok{mean}\NormalTok{((lstm\_preds\_original }\SpecialCharTok{{-}}\NormalTok{ true\_prices\_original)}\SpecialCharTok{\^{}}\DecValTok{2}\NormalTok{)}

\FunctionTok{cat}\NormalTok{(}\StringTok{"Test MSE (RNN, original scale):"}\NormalTok{, mse\_rnn, }\StringTok{"}\SpecialCharTok{\textbackslash{}n}\StringTok{"}\NormalTok{)}
\end{Highlighting}
\end{Shaded}

\begin{verbatim}
## Test MSE (RNN, original scale): 10.3685
\end{verbatim}

\begin{Shaded}
\begin{Highlighting}[]
\FunctionTok{cat}\NormalTok{(}\StringTok{"Test MSE (LSTM, original scale):"}\NormalTok{, mse\_lstm, }\StringTok{"}\SpecialCharTok{\textbackslash{}n}\StringTok{"}\NormalTok{)}
\end{Highlighting}
\end{Shaded}

\begin{verbatim}
## Test MSE (LSTM, original scale): 13.48961
\end{verbatim}

Plot 10-Day Forecast (Original Scale)

\begin{Shaded}
\begin{Highlighting}[]
\FunctionTok{plot}\NormalTok{(}
\NormalTok{  (}\FunctionTok{length}\NormalTok{(train\_data)}\SpecialCharTok{+}\DecValTok{1}\NormalTok{)}\SpecialCharTok{:}\NormalTok{(}\FunctionTok{length}\NormalTok{(train\_data)}\SpecialCharTok{+}\DecValTok{10}\NormalTok{),}
\NormalTok{  true\_prices\_original,}
  \AttributeTok{type =} \StringTok{\textquotesingle{}l\textquotesingle{}}\NormalTok{, }
  \AttributeTok{lwd =} \DecValTok{2}\NormalTok{, }
  \AttributeTok{col =} \StringTok{\textquotesingle{}black\textquotesingle{}}\NormalTok{, }
  \AttributeTok{xlab =} \StringTok{"Time Index"}\NormalTok{, }
  \AttributeTok{ylab =} \StringTok{"Close Price (USD)"}\NormalTok{,}
  \AttributeTok{main =} \StringTok{"10{-}Day Forecast vs True (Original Scale)"}
\NormalTok{)}

\FunctionTok{lines}\NormalTok{(}
\NormalTok{  (}\FunctionTok{length}\NormalTok{(train\_data)}\SpecialCharTok{+}\DecValTok{1}\NormalTok{)}\SpecialCharTok{:}\NormalTok{(}\FunctionTok{length}\NormalTok{(train\_data)}\SpecialCharTok{+}\DecValTok{10}\NormalTok{),}
\NormalTok{  rnn\_preds\_original, }
  \AttributeTok{col =} \StringTok{\textquotesingle{}red\textquotesingle{}}\NormalTok{, }
  \AttributeTok{lwd =} \DecValTok{2}
\NormalTok{)}

\FunctionTok{lines}\NormalTok{(}
\NormalTok{  (}\FunctionTok{length}\NormalTok{(train\_data)}\SpecialCharTok{+}\DecValTok{1}\NormalTok{)}\SpecialCharTok{:}\NormalTok{(}\FunctionTok{length}\NormalTok{(train\_data)}\SpecialCharTok{+}\DecValTok{10}\NormalTok{),}
\NormalTok{  lstm\_preds\_original, }
  \AttributeTok{col =} \StringTok{\textquotesingle{}blue\textquotesingle{}}\NormalTok{, }
  \AttributeTok{lwd =} \DecValTok{2}
\NormalTok{)}

\FunctionTok{legend}\NormalTok{(}
  \StringTok{"bottomleft"}\NormalTok{, }
  \AttributeTok{legend =} \FunctionTok{c}\NormalTok{(}\StringTok{"True Close"}\NormalTok{, }\StringTok{"RNN Prediction"}\NormalTok{, }\StringTok{"LSTM Prediction"}\NormalTok{),}
  \AttributeTok{col =} \FunctionTok{c}\NormalTok{(}\StringTok{"black"}\NormalTok{, }\StringTok{"red"}\NormalTok{, }\StringTok{"blue"}\NormalTok{), }
  \AttributeTok{lwd =} \DecValTok{2}
\NormalTok{)}
\end{Highlighting}
\end{Shaded}

\includegraphics{RNN-adn-LSTM_files/figure-latex/unnamed-chunk-41-1.pdf}

\hfill\break
Plot 1-Step Ahead Forecast Only (Original Scale)

\begin{Shaded}
\begin{Highlighting}[]
\FunctionTok{plot}\NormalTok{(}
  \DecValTok{1}\NormalTok{, }
\NormalTok{  true\_prices\_original[}\DecValTok{1}\NormalTok{], }
  \AttributeTok{type =} \StringTok{\textquotesingle{}p\textquotesingle{}}\NormalTok{, }
  \AttributeTok{lwd =} \DecValTok{2}\NormalTok{, }
  \AttributeTok{col =} \StringTok{\textquotesingle{}black\textquotesingle{}}\NormalTok{,}
  \AttributeTok{xlab =} \StringTok{"Forecast Horizon (1 Day)"}\NormalTok{, }
  \AttributeTok{ylab =} \StringTok{"Close Price (USD)"}\NormalTok{,}
  \AttributeTok{main =} \StringTok{"1{-}Day Forecast in Original Scale"}\NormalTok{,}
  \AttributeTok{xlim =} \FunctionTok{c}\NormalTok{(}\FloatTok{0.5}\NormalTok{, }\FloatTok{1.5}\NormalTok{),}
  \AttributeTok{ylim =} \FunctionTok{range}\NormalTok{(}\FunctionTok{c}\NormalTok{(true\_prices\_original[}\DecValTok{1}\NormalTok{], rnn\_preds\_original[}\DecValTok{1}\NormalTok{], lstm\_preds\_original[}\DecValTok{1}\NormalTok{]))}
\NormalTok{)}

\FunctionTok{axis}\NormalTok{(}
  \DecValTok{1}\NormalTok{, }
  \AttributeTok{at =} \DecValTok{1}\NormalTok{, }
  \AttributeTok{labels =} \FunctionTok{paste0}\NormalTok{(}\StringTok{"Day "}\NormalTok{, }\FunctionTok{length}\NormalTok{(train\_data)}\SpecialCharTok{+}\DecValTok{1}\NormalTok{)}
\NormalTok{)}

\FunctionTok{points}\NormalTok{(}
  \DecValTok{1}\NormalTok{, }
\NormalTok{  rnn\_preds\_original[}\DecValTok{1}\NormalTok{], }
  \AttributeTok{col =} \StringTok{\textquotesingle{}red\textquotesingle{}}\NormalTok{, }
  \AttributeTok{pch =} \DecValTok{17}
\NormalTok{)}

\FunctionTok{points}\NormalTok{(}
  \DecValTok{1}\NormalTok{, }
\NormalTok{  lstm\_preds\_original[}\DecValTok{1}\NormalTok{], }
  \AttributeTok{col =} \StringTok{\textquotesingle{}blue\textquotesingle{}}\NormalTok{, }
  \AttributeTok{pch =} \DecValTok{17}
\NormalTok{)}

\FunctionTok{legend}\NormalTok{(}
  \StringTok{"topleft"}\NormalTok{, }
  \AttributeTok{legend =} \FunctionTok{c}\NormalTok{(}\StringTok{"True Close"}\NormalTok{, }\StringTok{"RNN Forecast"}\NormalTok{, }\StringTok{"LSTM Forecast"}\NormalTok{),}
  \AttributeTok{col =} \FunctionTok{c}\NormalTok{(}\StringTok{"black"}\NormalTok{, }\StringTok{"red"}\NormalTok{, }\StringTok{"blue"}\NormalTok{), }
  \AttributeTok{pch =} \FunctionTok{c}\NormalTok{(}\DecValTok{16}\NormalTok{, }\DecValTok{17}\NormalTok{, }\DecValTok{17}\NormalTok{)}
\NormalTok{)}
\end{Highlighting}
\end{Shaded}

\includegraphics{RNN-adn-LSTM_files/figure-latex/unnamed-chunk-42-1.pdf}

Model Test MSE (Original USD Scale)

\begin{itemize}
\item
  RNN 13.13
\item
  LSTM 10.63

  LSTM has lower MSE than RNN. LSTM outperforms RNN when predicting
  actual stock prices.
\end{itemize}

\end{document}
